\documentclass[12pt]{article}
\usepackage{fullpage}
\usepackage{amsmath,amssymb,amsthm}

\newtheorem{theorem}{Theorem}
\newtheorem{lemma}{Lemma}
\newtheorem{claim}{Claim}

\newcommand{\Aut}{\operatorname{Aut}}
\newcommand{\spanv}{\operatorname{span}}

\begin{document}

\begin{center}
{\bf A computably AUT-countable example with no finite existential
automorphism basis}
\end{center}

\begin{theorem}
There is a computable countable structure $\mathcal A$ such that every
automorphism of every copy of $\mathcal A$ is computable from that copy,
but for every finite parameter set $P\subseteq A$ there is an automorphism
$\pi$ of $\mathcal A$ whose graph is not $\Sigma^{in}_1$-definable over
$P\cup\pi(P)$.
\end{theorem}

\begin{proof}
Let $V=[\omega]^{<\omega}$, with addition $+$ given by symmetric
difference; thus $V$ is a countably infinite vector space over
$\mathbb F_2$.  Fix a computable sequence $(D_i)_{i\in\omega}$ of finite
subsets of $V$ such that every finite subset of $V$ occurs infinitely
often, indeed cofinally often.

The language has two unary predicates $I,V$, a binary relation $<$ on
$I$, and a ternary relation $R$.  The structure $\mathcal A$ has two
sorts: $I^\mathcal A=\omega$ with its usual order, and
$V^\mathcal A=V$.  For $i\in I$ and $x,y\in V$, put
\[
        R(i,x,y)\quad\Longleftrightarrow\quad x+y\in D_i,
\]
and let $R$ be false on all other triples.  This is a computable
relational structure.

First compute its automorphism group.  Since $(I,<)$ has order type
$\omega$, every automorphism fixes $I$ pointwise.  If $f\in\Aut(\mathcal
A)$, let $c=f(0_V)$.  For every $h\in V$ choose $i$ with $D_i=\{h\}$.
Then $R(i,x,x+h)$ for all $x\in V$, and preservation of $R$ gives
$f(x+h)=f(x)+h$.  Taking $x=0_V$ and then writing an arbitrary element as
$h$, we get
\[
        f(h)=c+h  \qquad(h\in V).
\]
Conversely, for each $c\in V$, the map $\tau_c$ which fixes $I$
pointwise and sends $x\in V$ to $x+c$ preserves all differences $x+y$,
hence preserves $R$.  Thus
\[
        \Aut(\mathcal A)=\{\tau_c:c\in V\}.
\]

This already implies computable AUT-countability, with no cone oracle.
Let $\mathcal B\cong\mathcal A$, and let $\sigma\in\Aut(\mathcal B)$.
Transporting $\sigma$ to $\mathcal A$, it is some translation $\tau_c$.
Choose an element $i$ of the $I$-sort of $\mathcal B$ corresponding to
some index with $D_i=\{c\}$.  A Turing program with oracle the atomic
diagram of $\mathcal B$ and with this natural number $i$ hardwired
computes $\sigma$: it outputs $n$ on inputs $n$ in the $I$-sort, and on
inputs $x$ in the $V$-sort it searches for the unique $y$ with
$R(i,x,y)$.  Hence every automorphism of every copy is computable
relative to the copy.

It remains to prove the promised failure of $\Sigma^{in}_1$-definability.
Fix a finite parameter set $P\subseteq A$.  Let
\[
        M=\max(P\cap I)
\]
if $P\cap I\neq\varnothing$, and put $M=-1$ otherwise.  Let
\[
        W=\spanv\Bigl(\bigcup_{i\le M}D_i\Bigr)\le V
\]
(where $W=\{0\}$ if $M=-1$).  Choose a nonzero $h\in V\setminus W$, and
let $\pi=\tau_h$.  Put $Q=P\cup\pi(P)$.  Since $W$ is finite-dimensional,
choose $x\in V$ such that the two cosets
\[
        x+W,\qquad x+h+W
\]
are disjoint from the finite set $Q\cap V$.  Let $y=x+h$.

We shall show that the graph-pair $(x,y)$ is not isolated inside
$\operatorname{graph}(\pi)$ by any existential formula over $Q$.  This
suffices, since a $\Sigma^{in}_1$ definition would contain some
existential disjunct true of $(x,y)$ and contained in the graph.

Let
\[
        \exists \bar z\,\theta(\bar q,u,v,\bar z)
\]
be an existential formula with parameters $\bar q$ from $Q$, and suppose
it is true at $(u,v)=(x,y)$.  Choose witnesses in $\mathcal A$.  We may
work with one true conjunction of literals occurring in a disjunctive
normal form of the quantifier-free formula $\theta$; preserving the truth
values of the finitely many atomic formulas appearing in that conjunction
will preserve truth of the original formula.

Call a marker from the $I$-sort \emph{low} if its value is at most $M$.
All low markers have their associated finite set $D_i$ contained in
$W$.  Let $F$ be the finite set of all $V$-sort elements occurring among
$x,y$, the parameters, and the chosen witnesses.  We now choose
$t\in V$ and change precisely the elements of $F$ lying in the coset
$C=y+W$ by adding $t$.  The coset $C$ contains no parameter and does not
contain $x$.

There are only finitely many forbidden choices of $t$.  We avoid
$t=0$ and $t=h$, avoid choices causing a changed vector to collide with
an unchanged one, and, for every low marker $i\le M$ and every pair
$a\in C$, $b\in F\setminus C$, avoid
\[
        a+b+t\in D_i.
\]
These exclusions preserve all equalities and all $R$-literals whose
marker is low: within $C$ all differences are unchanged; outside $C$
nothing is changed; and across the cut no new difference falls in any
low $D_i$.

It remains to deal with marker variables whose chosen values are above
$M$.  They are not bounded by the parameters in the order, so they may be
replaced by sufficiently large markers.  After the above vector shift,
for each equality-class $\mu$ of such high marker variables, collect the
finite set $A_\mu$ of differences required by positive $R$-literals with
marker $\mu$, and the finite set $B_\mu$ of differences forbidden by
negative $R$-literals with marker $\mu$.  We also require
$A_\mu\cap B_\mu=\varnothing$.  This only excludes finitely many further
values of $t$: an equality between one shifted positive difference and
one shifted negative difference is either impossible already in the
original true diagram, or else is a single linear equation for $t$.
Choose $t$ avoiding all these finitely many obstructions.  Since every
finite subset of $V$ occurs cofinally often among the $D_i$'s, we can
choose new indices above $M$, in the same order and equality pattern as
the original high marker witnesses, with associated sets exactly
$A_\mu$.  Then every displayed $R$-literal with a high marker has the
same truth value as before.  Order and equality literals among markers
are preserved by construction.

Thus the same existential formula is true of the new pair
\[
        (x,\;y+t)=(x,\;x+h+t).
\]
But $t\neq 0$, so $y+t\neq x+h=\pi(x)$.  Hence this new pair is not in
the graph of $\pi$.  Therefore no existential formula over $Q$ which is
true of $(x,\pi(x))$ can be contained in the graph of $\pi$, and the graph
of $\pi$ is not $\Sigma^{in}_1$-definable over $Q=P\cup\pi(P)$.

Since $P$ was arbitrary, the required property holds for every finite
parameter set.
\end{proof}

\end{document}
